\documentclass{article}

% Packages
\usepackage{amsmath, amsfonts, amssymb, amsthm, bm, bbm, comment} 
\usepackage{enumerate, graphicx, xcolor, tcolorbox, hyphenat, circledsteps}
\usepackage[shortlabels]{enumitem}
\usepackage{url, hyperref}
\hypersetup{
	colorlinks   = true,
	urlcolor     = cyan,
	linkcolor    = red,
	citecolor    = black
}

% Formatting
\usepackage[a4paper, noheadfoot, scale={0.7,0.8}]{geometry}
\title{Wk 7 Supervision Problems 1}
\author{Luca Bollini}
\date{17 November 2026}
\hyphenpenalty=10000

% Commands
\newcommand{\N}{\mathbb{N}}
\newcommand{\Z}{\mathbb{Z}}
\newcommand{\Q}{\mathbb{Q}}
\newcommand{\R}{\mathbb{R}}
\newcommand{\eps}{\varepsilon}
\newcommand{\ddt}[1]{\frac{d}{d#1}}
\newcommand{\ddttwo}[1]{\frac{d^2}{d#1^2}}
\newcommand{\ts}{\thinspace}

% Colours
\newtcolorbox{questionbox}{
	colframe=black, 
	colback=white, 
	boxrule=0.2mm,
	arc=0mm,
}

\begin{document}
\maketitle

\begin{center}
\huge{\textbf{Cauchy Sequences}}
\end{center}

\begin{comment}
    Students struggled with the first question. Remind them that the Cauchy property is much stronger than just controlling the difference between consequtive terms, hence why the r^n factor is needed for part i).

    Part i) is a good exercise in using the triangle inequality. Have them write up proper answers in sentences, and let them grade each other's work. Use this as a way of illustrating good clarity. 
\end{comment}

\begin{enumerate}
    \item Let $(s_n)$ and $(t_n)$ be sequences of real numbers, and suppose that $(s_n)$ converges. Does $(t_n)$ converge if
    \begin{enumerate}[(i)]
        \item there exists $0<r<1$ such that for all $n$ sufficiently large,
        \[
        |t_{n+1}-t_{n}| \leq r^n|s_{n+1}-s_n|\text{?}
        \]
        \item for any $n\in\N$
        \[
        |t_{n+1}-t_{n}| \leq |s_{n+1}-s_n|\text{?}
        \]
        \item there exists $N\in\N$ such that for all $m,n>N$,
        \[
        |t_{n+m}-t_{n}| \leq |s_{n+m}-s_n|\text{?}
        \]
    \end{enumerate}

    \item Define the sequence $(a_n)$ by
    \[
    a_n:= 1+\frac{1}{2!}+\frac{1}{3!}+\cdots+\frac{1}{n!}.
    \]
    Using the Cauchy property of sequences, prove that $(a_n)$ converges.
\end{enumerate}

\begin{center}
\huge{\textbf{Nondimensionalisation}}
\end{center}

\begin{enumerate}
    \item Consider the logistic equation
    \[
    \dot N = rN\left(1-\frac NK\right)
    \]
    subject to $N(0) = N_0$.
    \begin{enumerate}[(i)]
        \item Find the dimensions of the parameters $r,K$ and $N_0$, in terms of the dimensions of $N$.
        \item Show the system can be rewritten with the dimensionless formula
        \[
        \frac{dx}{d\tau} = x(1-x),\quad x(0) = x_0
        \]
        for an appropriate choice of dimensionless variables $x,x_0$, and  $\tau$.
        \item Find a different nondimensionalisation in terms of variables $u$ and $\tau$, chosen such that the initial condition is always $u_0=1$.
        \item Can you think of any advantages of one nondimensionalisation over the other?
    \end{enumerate}
\end{enumerate}

\end{document}